% Mathématiques 4e — cours V3
% Triangles, Thalès et cosinus

\section{Objectifs}

\begin{itemize}
\item Utiliser des cas d'égalité de triangles dans des raisonnements simples.
\item Reconnaître une configuration de Thalès avec triangles emboîtés.
\item Utiliser le théorème de Thalès pour calculer une longueur.
\item Utiliser la réciproque du théorème de Thalès pour prouver un parallélisme.
\item Définir le cosinus d'un angle aigu dans un triangle rectangle.
\item Calculer une longueur avec le cosinus.
\item Déterminer un angle à l'aide de la calculatrice.
\item Rédiger une démonstration structurée.
\item Distinguer Thalès, Pythagore et cosinus selon les données disponibles.
\end{itemize}

\section{1. Triangles égaux}

Deux triangles sont égaux lorsqu'ils peuvent être superposés exactement.

Ils ont alors :
\begin{itemize}
\item des côtés correspondants de même longueur ;
\item des angles correspondants de même mesure.
\end{itemize}

Cette idée permet de justifier des égalités géométriques sans mesurer sur la figure.

\section{2. Cas simple côté-côté-côté}

Si deux triangles ont leurs trois côtés respectivement de mêmes longueurs, ils sont égaux.

Par exemple :
\[
AB=DE,
\]
\[
BC=EF,
\]
\[
CA=FD.
\]

On peut alors déduire l'égalité des angles correspondants.

\section{3. Cas côté-angle-côté}

Si deux triangles possèdent :
\begin{itemize}
\item deux côtés respectivement égaux ;
\item l'angle compris entre ces côtés de même mesure,
\end{itemize}
alors ils sont égaux.

Cette propriété peut intervenir dans une démonstration de symétrie ou de construction.

\section{4. Configuration de Thalès}

On considère un triangle $ABC$.

Les points $M$ et $N$ appartiennent respectivement aux côtés $[AB]$ et $[AC]$.

Si :
\[
(MN)\parallel(BC),
\]
alors les triangles :
\[
AMN
\]
et :
\[
ABC
\]
ont la même forme et des longueurs proportionnelles.

\coursefigure{/images/mathematiques/4eme/thales-configuration.svg}{Triangle ABC contenant le segment MN parallèle au côté BC.}{Dans cette configuration emboîtée, les longueurs correspondantes des deux triangles sont proportionnelles.}

\section{5. Théorème de Thalès}

\begin{propriete}
Si :
\[
M\in[AB],
\quad
N\in[AC],
\quad
(MN)\parallel(BC),
\]
alors :
\[
\boxed{
\dfrac{AM}{AB}
=
\dfrac{AN}{AC}
=
\dfrac{MN}{BC}
}.
\]
\end{propriete}

Les trois rapports utilisent des longueurs correspondantes dans le même ordre.

\section{6. Calculer une longueur}

Supposons :
\[
AM=3,
\quad
AB=5,
\quad
AC=8,
\]
et :
\[
(MN)\parallel(BC).
\]

Par Thalès :
\[
\dfrac{AM}{AB}=\dfrac{AN}{AC}.
\]

Donc :
\[
\dfrac35=\dfrac{AN}{8}.
\]

Produit en croix :
\[
5AN=24.
\]

Ainsi :
\[
\boxed{AN=4{,}8}.
\]

\section{7. Autre longueur}

Avec :
\[
AM=4,
\quad
AB=10,
\quad
BC=15,
\]
on cherche $MN$.

\[
\dfrac4{10}=\dfrac{MN}{15}.
\]

Donc :
\[
10MN=60.
\]

Ainsi :
\[
\boxed{MN=6}.
\]

\section{8. Bien associer les côtés}

Une erreur fréquente consiste à mélanger les côtés correspondants.

Le petit triangle est :
\[
AMN.
\]

Le grand triangle est :
\[
ABC.
\]

Les correspondances sont :
\[
AM\leftrightarrow AB,
\]
\[
AN\leftrightarrow AC,
\]
\[
MN\leftrightarrow BC.
\]

\section{9. Réciproque de Thalès}

\begin{propriete}
Dans une configuration où les points sont correctement alignés, si les rapports de longueurs correspondantes sont égaux, alors on peut conclure au parallélisme.
\end{propriete}

Par exemple, si :
\[
\dfrac{AM}{AB}
=
\dfrac{AN}{AC},
\]
avec $M$ sur $[AB]$ et $N$ sur $[AC]$, alors :
\[
(MN)\parallel(BC).
\]

\section{10. Exemple de réciproque}

On connaît :
\[
AM=4,
\quad
AB=10,
\quad
AN=6,
\quad
AC=15.
\]

On calcule :
\[
\dfrac{AM}{AB}=\dfrac4{10}=0{,}4,
\]
et :
\[
\dfrac{AN}{AC}=\dfrac6{15}=0{,}4.
\]

Les rapports sont égaux.

Sous les conditions d'alignement adaptées :
\[
\boxed{(MN)\parallel(BC)}.
\]

\section{11. Thalès et agrandissement}

Le rapport :
\[
\dfrac{AB}{AM}
\]
est un coefficient d'agrandissement entre les deux triangles.

Si :
\[
\dfrac{AB}{AM}=2{,}5,
\]
toutes les longueurs du grand triangle sont :
\[
2{,}5
\]
fois celles du petit triangle.

\section{12. Triangle rectangle et cosinus}

Dans un triangle rectangle, un angle aigu permet de distinguer :
\begin{itemize}
\item l'hypoténuse ;
\item le côté adjacent à cet angle ;
\item le côté opposé.
\end{itemize}

Le cosinus compare le côté adjacent à l'hypoténuse.

\coursefigure{/images/mathematiques/4eme/cosinus-triangle.svg}{Triangle rectangle avec angle alpha, côté adjacent et hypoténuse identifiés.}{Pour un angle aigu dans un triangle rectangle, le cosinus est le rapport longueur du côté adjacent sur longueur de l'hypoténuse.}

\section{13. Définition du cosinus}

\begin{definition}
Dans un triangle rectangle :
\[
\boxed{
\cos(\alpha)
=
\dfrac{\text{longueur du côté adjacent à }\alpha}
{\text{longueur de l'hypoténuse}}
}.
\]
\end{definition}

Comme le côté adjacent est plus court que l'hypoténuse :
\[
0<\cos(\alpha)<1
\]
pour un angle aigu.

\section{14. Calculer un côté adjacent}

Dans un triangle rectangle, l'hypoténuse mesure :
\[
10\ \text{cm}
\]
et :
\[
\alpha=60^\circ.
\]

On sait :
\[
\cos60^\circ=0{,}5.
\]

Donc :
\[
\dfrac{\text{adjacent}}{10}=0{,}5.
\]

Ainsi :
\[
\text{adjacent}=5\ \text{cm}.
\]

\section{15. Calculer l'hypoténuse}

Le côté adjacent à un angle de :
\[
40^\circ
\]
mesure :
\[
7\ \text{cm}.
\]

Si $h$ est l'hypoténuse :
\[
\cos40^\circ=\dfrac7h.
\]

Donc :
\[
h=\dfrac7{\cos40^\circ}.
\]

À la calculatrice :
\[
h\approx9{,}14\ \text{cm}.
\]

\section{16. Calculer un angle}

Dans un triangle rectangle :
\[
\text{adjacent}=6,
\]
\[
\text{hypoténuse}=10.
\]

Donc :
\[
\cos(\alpha)=\dfrac6{10}=0{,}6.
\]

On utilise la fonction inverse de la calculatrice :
\[
\alpha\approx53{,}1^\circ.
\]

\section{17. Mode degré}

Pour des angles exprimés en degrés, la calculatrice doit être réglée en mode degré.

Une erreur de mode peut donner une valeur totalement incohérente.

\begin{attention}
Toujours contrôler qu'un angle aigu trouvé vérifie :
\[
0^\circ<\alpha<90^\circ.
\]
\end{attention}

\section{18. Cosinus ou Pythagore ?}

Pythagore est adapté lorsque deux longueurs d'un triangle rectangle sont connues et que l'on cherche la troisième.

Le cosinus est adapté lorsqu'on connaît :
\begin{itemize}
\item un angle aigu et une longueur ;
\item ou deux longueurs impliquant le côté adjacent et l'hypoténuse.
\end{itemize}

\section{19. Thalès ou cosinus ?}

Thalès utilise un parallélisme et des triangles de même forme.

Le cosinus utilise un triangle rectangle et un angle aigu.

Avant tout calcul, il faut identifier la structure géométrique.

\section{20. Problème de pente}

Une rampe de longueur :
\[
5\ \text{m}
\]
forme un angle :
\[
12^\circ
\]
avec le sol.

Sa longueur horizontale $d$ vérifie :
\[
\cos12^\circ=\dfrac d5.
\]

Donc :
\[
d=5\cos12^\circ.
\]

Ainsi :
\[
d\approx4{,}89\ \text{m}.
\]

\section{21. Hauteur indirecte}

Un câble de :
\[
12\ \text{m}
\]
forme un angle de :
\[
35^\circ
\]
avec le sol.

Le cosinus permet d'obtenir la composante horizontale :
\[
d=12\cos35^\circ.
\]

Pour calculer directement une hauteur opposée à l'angle, le cosinus seul n'est pas forcément le rapport le plus immédiat ; on peut compléter la figure ou utiliser Pythagore après avoir trouvé $d$.

\section{22. Démonstration structurée}

Une démonstration doit distinguer :
\begin{itemize}
\item les données ;
\item la propriété utilisée ;
\item les calculs ;
\item la conclusion.
\end{itemize}

Une figure n'est pas une preuve.

\section{23. Méthode Thalès}

\begin{methode}
\begin{enumerate}
\item identifier les deux triangles ;
\item vérifier les alignements et le parallélisme ;
\item écrire les côtés correspondants dans le même ordre ;
\item écrire l'égalité des rapports ;
\item remplacer par les valeurs ;
\item calculer la longueur cherchée ;
\item donner l'unité.
\end{enumerate}
\end{methode}

\section{24. Méthode cosinus}

\begin{methode}
\begin{enumerate}
\item vérifier que le triangle est rectangle ;
\item repérer l'angle étudié ;
\item identifier hypoténuse et côté adjacent ;
\item écrire la définition du cosinus ;
\item remplacer par les valeurs ;
\item isoler l'inconnue ;
\item utiliser la calculatrice si nécessaire ;
\item contrôler l'ordre de grandeur.
\end{enumerate}
\end{methode}

\section{25. Erreurs fréquentes}

\begin{itemize}
\item Appliquer Thalès sans parallélisme.
\item Croiser les correspondances entre petit et grand triangle.
\item Utiliser la réciproque sans vérifier les alignements.
\item Confondre côté adjacent et côté opposé.
\item Choisir un côté autre que l'hypoténuse au dénominateur du cosinus.
\item Utiliser la calculatrice en mode radian.
\item Mesurer sur la figure au lieu d'utiliser les propriétés.
\end{itemize}

\section{À retenir}

\begin{aretenir}
Dans une configuration de Thalès :
\[
\boxed{
\dfrac{AM}{AB}
=
\dfrac{AN}{AC}
=
\dfrac{MN}{BC}
}.
\]

L'égalité des rapports peut aussi permettre de prouver un parallélisme par réciproque.

Dans un triangle rectangle :
\[
\boxed{
\cos(\alpha)
=
\dfrac{\text{côté adjacent}}{\text{hypoténuse}}
}.
\]

Choisir entre Pythagore, Thalès et cosinus dépend toujours des données géométriques disponibles.
\end{aretenir}
